% answers_short.tex —— 2006 年·普通高等学校招生考试·湖南卷理科 · 简短答案速查
% 由 tools/make_short_answers.py 自 answers.tex 的 \gksolhead/\gkanswer 提取，
% 并与 2006/2006_数学真题_湖南理_答案要点.md 逐题比对后生成，勿手改（改源 MD 后重跑）。

\begin{gknotice}
\textbf{使用说明}：本节只给最终答案便于快速核对，完整推导见解析版。答案为编写者独立推导，\textbf{非官方标准答案}。
\end{gknotice}

\gkpart{选择题}
\begin{center}\begin{tabular}{ccccccccccc}
\toprule
题号 & 1 & 2 & 3 & 4 & 5 & 6 & 7 & 8 & 9 & 10 \\
\midrule
答案 & D & A & D & A & B & D & A & C & C & B \\
\bottomrule
\end{tabular}\end{center}

\gkpart{填空题}
\begin{itemize}[leftmargin=2.6em]
  \item[\textbf{11.}] $-2$。
  \item[\textbf{12.}] $5$。
  \item[\textbf{13.}] $\dfrac34$。
  \item[\textbf{14.}] 例如 $(a,b)=(1,-1)$。
  \item[\textbf{15.}] $x\in(-\infty,0)$；当 $x=-\dfrac12$ 时，$y\in\left(\dfrac12,\dfrac32\right)$。
\end{itemize}

\gkpart{解答题}
\begin{itemize}[leftmargin=2.6em]
  \item[\textbf{16.}] (1) 证明见解析；(2) $\beta=\dfrac\pi3$
  \item[\textbf{17.}] (1) $0.31$；(2) $2.5$ 家；(3) U/PENDING
  \item[\textbf{18.}] (1) $PQ\perp$ 平面 $ABCD$；(2) $\arccos\dfrac1{3\sqrt3}$；(3) $\dfrac3{\sqrt2}$
  \item[\textbf{19.}] (1) $0<a_{n+1}<a_n<1$；(2) $a_{n+1}<\dfrac16a_n^3$
  \item[\textbf{20.}] (1) 方案甲用水 $19$ 单位；当 $c=0.95$ 时，方案乙初洗用水 $3$ 单位、二次用水 $4a$ 单位，总用水 $3+4a$ 单位，少于方案甲；(2) 令 $c=1-\dfrac1{\sqrt{500a}}$；初洗用水 $x=\sqrt{20a}-1$，二次用水 $y=\sqrt{20a}-a$，最少总用水为 $2\sqrt{20a}-a-1$。该最少值在 $1\le a\le3$ 上随 $a$ 严格增加
  \item[\textbf{21.}] (1) $m=0,\ p=\dfrac98$；抛物线焦点为 $\left(\dfrac9{16},0\right)$，不在 $AB:x=1$ 上；(2) 存在：$m=\pm\dfrac{\sqrt6}3$、$p=\dfrac43$；公共弦为 $y=-3m(x-1)$
\end{itemize}
